% ─── Capítulo 8: Conclusiones ─────────────────────────────────────────────────
\chapter{Conclusiones y Trabajo Futuro}

\section{Resumen del Proceso de Compilación}

El compilador TPP implementa el pipeline completo de traducción desde un lenguaje de programación imperativo en español hasta instrucciones ejecutables en hardware real. A continuación se resume el recorrido de un programa a través de cada fase:

\begin{table}[H]
\centering
\begin{tabular}{p{2.5cm}p{3.5cm}p{7cm}}
\toprule
\textbf{Fase} & \textbf{Archivo} & \textbf{Transformación} \\
\midrule
Léxico & \texttt{lexer.l} & Texto → tokens (\texttt{NUM\_ENTERO}, \texttt{ID}, \texttt{SI}, ...) \\
Sintáctico & \texttt{parser.y} & Tokens → AST jerárquico \\
Semántico & \texttt{semantic.c} & AST → AST anotado con tipos y offsets de memoria \\
Tabla de símb. & \texttt{symtab.c} & Registro de variables/funciones con sus direcciones \\
Optimización & \texttt{opt.c} & AST → AST simplificado (plegado, algebraica, muerto) \\
IR/TAC & \texttt{ir.c} & AST → representación plana de 3 direcciones (diagnóstico) \\
Codegen & \texttt{codegen.c} & AST → instrucciones RV32I en \texttt{output.s} \\
\bottomrule
\end{tabular}
\caption{Resumen del pipeline completo}
\label{tab:resumen}
\end{table}

\section{Logros Implementados}

\subsection{Frontend}
\begin{itemize}
    \item \textbf{Analizador léxico} completo con 30+ tokens, soporte de comentarios, cadenas con secuencias de escape, y reporte preciso de errores con línea y columna.
    \item \textbf{Analizador sintáctico} LALR(1) con gramática completa, resolución de ambigüedades (precedencia de operadores, dangling else), y recuperación de errores en modo pánico.
    \item \textbf{Analizador semántico} con verificación de tipos, detección de variables no declaradas, redefinición de símbolos y anotación de direcciones de memoria en el AST.
    \item \textbf{Tabla de símbolos} con scoping anidado, asignación estática de memoria y visualización detallada al final del análisis.
\end{itemize}

\subsection{Backend}
\begin{itemize}
    \item \textbf{Optimizador de AST} con 4 pases: plegado de constantes, simplificación algebraica, eliminación de código muerto y propagación de constantes.
    \item \textbf{Generador de código intermedio} (TAC) para diagnóstico y comprensión del proceso.
    \item \textbf{Generador de código RV32I} completo con:
    \begin{itemize}
        \item Generación bifásica para soporte de funciones.
        \item Convención de llamada RISC-V (a0--a7, ra).
        \item Prólogo y epílogo completo de funciones (salvar/restaurar ra).
        \item Acceso a arreglos con cálculo dinámico de dirección.
        \item Multiplicación por software (\texttt{\_\_soft\_mul}).
        \item Impresión de enteros como ASCII (\texttt{\_\_print\_num}).
        \item Salida TTY a la dirección Logisim \texttt{0xff000004}.
    \end{itemize}
\end{itemize}

\section{Construcciones Soportadas}

\begin{table}[H]
\centering
\begin{tabular}{ll}
\toprule
\textbf{Característica} & \textbf{Estado} \\
\midrule
Variables enteras y decimales & \checkmark\ Implementado \\
Arreglos estáticos & \checkmark\ Implementado \\
Operadores aritméticos (+, -, *, /) & \checkmark\ Implementado \\
Operadores relacionales (==, !=, <, >, <=, >=) & \checkmark\ Implementado \\
Estructura condicional si/sino & \checkmark\ Implementado \\
Bucle mientras & \checkmark\ Implementado \\
Bucle para & \checkmark\ Implementado \\
Estructura depende/caso & \checkmark\ Implementado \\
Declaración y llamada de funciones & \checkmark\ Implementado \\
Parámetros por valor & \checkmark\ Implementado \\
Variables globales & \checkmark\ Implementado \\
Imprimir cadenas y enteros & \checkmark\ Implementado \\
Comentarios de línea (//) & \checkmark\ Implementado \\
Optimización de constantes & \checkmark\ Implementado \\
Recursión & $\sim$\ Parcial (requiere stack dinámico) \\
Tipo decimal en codegen & $\times$\ Pendiente \\
Paso de arreglos a funciones & $\times$\ Pendiente \\
\bottomrule
\end{tabular}
\caption{Estado de las características del compilador}
\label{tab:estado}
\end{table}

\section{Limitaciones Actuales}

\subsection{Modelo de Memoria Estático}

El compilador TPP utiliza un modelo de memoria \textbf{completamente estático}: todos los offsets se calculan en tiempo de compilación y son fijos. Esto tiene las siguientes implicaciones:

\begin{itemize}
    \item \textbf{Recursión limitada}: Una función recursiva sobreescribiría sus propias variables locales al ser llamada de nuevo, ya que todas las invocaciones comparten el mismo espacio de memoria.
    \item \textbf{Sin allocación dinámica}: No existe instrucción \texttt{malloc} equivalente.
    \item \textbf{Arreglos de tamaño fijo}: Los arreglos deben declararse con tamaño constante conocido en compilación.
\end{itemize}

\subsection{Sin guión bajo en identificadores}

El lexer solo reconoce \texttt{[a-zA-Z][a-zA-Z0-9]*} como identificadores. El carácter \texttt{\_} produce un error léxico. Se recomienda camelCase.

\section{Trabajo Futuro}

Las siguientes mejoras podrían extender las capacidades del compilador:

\begin{enumerate}
    \item \textbf{Stack dinámico por función}: Implementar un frame pointer (\texttt{fp}) separado del stack pointer para que cada invocación de función tenga su propio frame. Esto habilitaría la recursión completa.
    
    \item \textbf{Soporte de tipo decimal en codegen}: Generar instrucciones de punto flotante o emular operaciones decimales con enteros de punto fijo.
    
    \item \textbf{Paso de arreglos por referencia}: Pasar la dirección base del arreglo como puntero al argumento \texttt{a0}, permitiendo que las funciones modifiquen arreglos externos.
    
    \item \textbf{Múltiples archivos fuente}: Soporte para \texttt{importar} otros archivos \texttt{.to} y compilarlos juntos.
    
    \item \textbf{Optimizaciones adicionales}: Eliminación de subexpresiones comunes, renombramiento de registros, inlining de funciones pequeñas.
    
    \item \textbf{Soporte de guión bajo}: Extender el lexer para permitir \texttt{\_} en identificadores (\texttt{[a-zA-Z\_][a-zA-Z0-9\_]*}).
    
    \item \textbf{Extensión de ISA}: Soporte para la extensión \texttt{M} de RISC-V (multiplicación nativa con \texttt{mul}, \texttt{div}) una vez que el procesador Logisim la implemente.
\end{enumerate}

\section{Ejemplo Completo: De Fuente a Binario}

El siguiente ejemplo muestra el recorrido completo de una línea de código:

\begin{lstlisting}[style=tppstyle, caption={Código fuente TPP}]
entero resultado = 10 * 2 + 3;
imprimir(resultado);
\end{lstlisting}

\textbf{Tokens producidos por el lexer:}
\begin{verbatim}
TIPO_ENTERO  ID("resultado")  ASIG  NUM_ENTERO(10)  MULT  
NUM_ENTERO(2)  MAS  NUM_ENTERO(3)  PUNTOCOMA
IMPRIMIR  PARI  ID("resultado")  PARD  PUNTOCOMA
\end{verbatim}

\textbf{AST (antes de optimización):}
\begin{verbatim}
[Declaracion Var] entero
  [Asignacion] resultado
    [Operacion] '+'
      [Operacion] '*'
        [Entero] 10
        [Entero] 2
      [Entero] 3
[Imprimir]
  [Variable] resultado
\end{verbatim}

\textbf{AST (después de optimización — plegado en cascada):}
\begin{verbatim}
[Declaracion Var] entero
  [Asignacion] resultado
    [Entero] 23        ← 10*2=20, 20+3=23
[Imprimir]
  [Variable] resultado
\end{verbatim}

\textbf{Código RV32I generado:}
\begin{lstlisting}[style=asmstyle, caption={Ensamblador RV32I generado}]
    li  t0, 23         # Constante plegada (era 10*2+3)
    sw  t0, 0(sp)      # resultado = 23
    lw  t0, 0(sp)      # Cargar resultado
    mv  a0, t0
    call __print_num   # Imprimir el entero
    li  t0, 10         # Salto de línea
    sw  t0, 0xff000004
\end{lstlisting}

Sin la optimización de plegado, se habrían generado ~12 instrucciones adicionales para la multiplicación y la suma. Con la optimización, se reduce a 2 instrucciones esenciales.
